\chapter{Conclusões e Trabalho Futuro}
\label{chapter:conclusions}

\begin{introduction}
Este capítulo sintetiza os resultados obtidos, confronta-os com os objetivos definidos no Capítulo~\ref{chapter:introduction}, identifica as limitações do protótipo desenvolvido e propõe direções para trabalho futuro.
\end{introduction}

\section{Conclusões}
\label{sec:conclusoes}

O projeto visou desenvolver uma mesa sísmica bi-axial de baixo custo, com recurso a impressão 3D e a componentes eletrónicos de uso geral, para apoiar a avaliação de risco sísmico em países com recursos limitados. Os resultados obtidos confirmam que esse objetivo foi cumprido.

Do ponto de vista mecânico, a arquitetura baseada numa plataforma única suportada por quatro colunas com rótulas mostrou-se adequada para o movimento bi-axial. Ao contrário de abordagens com guias lineares rígidas por eixo, a inclinação livre das colunas permite que a plataforma se desloque simultaneamente nos dois eixos sem introduzir constrangimentos cinemáticos internos. Todos os componentes estruturais foram fabricados por impressão 3D em PLA, o que confirma a viabilidade desta tecnologia para o fabrico de equipamento científico de baixo custo e fácil reprodução.

Do ponto de vista eletrónico, a combinação ESP32 + CNC Shield ARD-CNC-Kit2 + quatro drivers DRV8825 + quatro motores NEMA17-04 demonstrou ser adequada para o controlo síncrono de quatro motores de passo em simultâneo. A configuração de microstepping $1/2$, com corrente ajustada para 1,68\,A por fase (V\textsubscript{ref} = 0,84\,V), garante o funcionamento dentro das especificações dos motores. A velocidade máxima exigida pelo perfil sísmico de El~Centro --- 1\,600\,steps/s na componente NS e 1\,000\,steps/s na componente EW --- situa-se bem dentro da gama operacional do hardware, o que valida a compatibilidade entre o registo sísmico escolhido e os componentes selecionados.

Do ponto de vista de software, o pipeline Python de pré-processamento converte os dados de deslocamento do ficheiro \texttt{ElCentro.xlsx} em arrays de passos de motor, que são embebidos diretamente no firmware do ESP32. A execução a 50\,Hz, com janelas de 20\,ms sincronizadas por \texttt{millis()}, reproduz o perfil sísmico de El~Centro (1940), com telemetria em tempo real via porta série disponível para monitorização e validação.

O custo de construção do protótipo ficou em aproximadamente 153\,\euro{}, o que representa menos de 1,6\% do custo estimado da solução comercial bi-axial equivalente mais acessível (Quanser Shake Table II XY, $>$10\,000\,\euro{}) e fica muito abaixo do limite de 500\,\euro{} definido como requisito não funcional. O requisito de reprodutibilidade foi igualmente satisfeito: todos os ficheiros de design 3D, firmware e scripts de pré-processamento estão disponíveis para consulta e reutilização.

Em suma, o projeto demonstrou que é possível construir uma mesa sísmica bi-axial funcional, reprodutível e de baixo custo com tecnologias acessíveis e amplamente disponíveis. O sistema responde diretamente à necessidade identificada pelo Professor Vítor Silva no seu trabalho de avaliação de risco estrutural em países em desenvolvimento e contribui para tornar o ensaio sísmico experimental acessível em contextos com recursos limitados.


\section{Limitações Identificadas}
\label{sec:limitacoes}

Apesar dos resultados positivos, o protótipo apresenta um conjunto de limitações que importa documentar.

\begin{itemize}
    \item \textbf{Controlo em malha aberta:} A ausência de encoders ou sensores de posição nos motores impede a deteção e correção de erros de passo. Se um motor perder passos por excesso de carga ou velocidade, o erro acumula-se silenciosamente e o movimento da plataforma diverge do perfil sísmico pretendido sem que o sistema o detete.

    \item \textbf{Margem de torque reduzida:} O torque estimado necessário por motor (5,6\,kg.cm) excede o torque de retenção nominal do NEMA17-04 (4,4\,kg.cm) com massas de modelo próximas do limite máximo. Por este motivo, recomenda-se não exceder 1\,kg de massa do modelo de teste, de modo a evitar perda de passo.

    \item \textbf{Resolução temporal do \texttt{millis()}:} A função \texttt{millis()} tem resolução de 1\,ms, o que introduz um \textit{jitter} de temporização de até $\pm$1\,ms por janela de 20\,ms (5\% do período). Ao longo de uma sequência sísmica completa, este erro pode causar uma ligeira deriva temporal entre o perfil executado e o perfil teórico.

    \item \textbf{Ausência de fins de curso:} A CNC Shield disponibiliza seis entradas para \textit{endstops}, mas estas não foram integradas na versão atual do firmware. Sem fins de curso, um erro nos dados ou no firmware pode levar a plataforma até ao limite físico de deslocamento, com risco de dano mecânico.

    \item \textbf{Modo de operação único implementado:} O firmware atual suporta apenas o modo de reprodução do registo de El~Centro. Os modos de frequência constante e de controlo manual, definidos como requisitos funcionais no Capítulo~\ref{chapter:requirements}, ficaram por implementar.

    \item \textbf{Rigidez estrutural limitada:} A estrutura em PLA apresenta menor rigidez e resistência térmica comparativamente a soluções em alumínio maquinado. As tolerâncias dimensionais das peças impressas em 3D afetam a precisão do alinhamento dos fusos, o que pode introduzir folgas e vibração residual na plataforma.

    \item \textbf{Capítulo de validação incompleto:} Os testes experimentais de verificação de deslocamentos e a comparação quantitativa entre o perfil simulado e o perfil real da plataforma não foram concluídos, o que constitui uma lacuna na validação do sistema (Capítulo~\ref{chapter:validation}).
\end{itemize}


\section{Propostas de Melhoria e Trabalho Futuro}
\label{sec:trabalho-futuro}

As limitações identificadas na secção anterior definem diretamente as principais direções de trabalho futuro.

\begin{itemize}
    \item \textbf{Integração de encoders e controlo em malha fechada:} A adição de encoders incrementais nos veios dos motores permitiria detetar e corrigir perdas de passo em tempo real, com ganhos significativos na fidelidade de reprodução do perfil sísmico. Uma implementação mais avançada poderia usar um controlador PID no ESP32 para fechar a malha de posição em cada eixo.

    \item \textbf{Integração dos fins de curso da CNC Shield:} A utilização das entradas de \textit{endstop} já disponíveis na CNC Shield, com micro-interruptores nos limites físicos de cada eixo, eliminaria o risco de colisão destrutiva e permitiria uma rotina de homing para posicionamento absoluto no arranque.

    \item \textbf{Implementação dos modos de frequência constante e manual:} O modo de frequência constante, com amplitude e frequência configuráveis, permitiria ensaios de varrimento para determinação das frequências naturais dos modelos estruturais. O modo manual, acessível via interface série ou Wi-Fi, serviria para posicionar a plataforma livremente durante a colocação e remoção dos modelos.

    \item \textbf{Interface de controlo via Wi-Fi:} O ESP32 integra Wi-Fi 802.11 b/g/n, que pode ser aproveitado para criar uma interface web de controlo acessível a partir de qualquer dispositivo na mesma rede local, sem necessidade de ligação USB durante os ensaios.

    \item \textbf{Suporte a registos sísmicos genéricos:} A versão atual do firmware tem os arrays de El~Centro codificados estaticamente. Uma evolução natural seria suportar o carregamento de registos sísmicos arbitrários em formato CSV a partir de um cartão microSD ou via Wi-Fi, o que tornaria o sistema compatível com qualquer registo da base de dados PEER~\cite{PEERDatabase}.

    \item \textbf{Validação experimental completa:} A realização dos testes de verificação de deslocamentos descritos no Capítulo~\ref{chapter:validation} --- com acelerómetro na plataforma, câmara de alta cadência ou régua óptica --- permitiria quantificar o erro de reprodução do perfil sísmico e validar formalmente o sistema face aos requisitos definidos.

    \item \textbf{Melhoria da rigidez mecânica:} A substituição das peças em PLA por peças em alumínio fresado nos elementos de maior solicitação mecânica --- suportes dos motores e blocos de fixação das porcas --- reduziria as folgas e melhoraria a fidelidade de reprodução a frequências mais elevadas.

    \item \textbf{Motorização com maior torque:} A substituição dos motores NEMA17 por motores NEMA23, com torque de retenção típico de 10--15\,kg.cm, aumentaria a margem de segurança e abriria o sistema ao ensaio de modelos estruturais mais pesados, o que alargaria o campo de aplicação.
\end{itemize}
